%% appB --- Błędne odczytania chaosu
%%
%% GENERATED by tools/gen_traps.py from notes/traps/chNN.json. Edit the
%% chapter's JSON, not this file; `make traps` rewrites it and
%% `check_structure.py --traps` checks it against the catalogue.

\chapter{Błędne odczytania chaosu}
\label{app:misreadings}

\noindent
Każde błędne odczytanie, które książka stara się poprawić, w kolejności, w jakiej wywołują je rozdziały. Każde jest zapisane głosem kogoś, kto w nie wierzy, a po nim następuje to, co jest prawdą, i podrozdział, w którym książka pozwala ci najpierw w nie wpaść. Jeśli trafiłeś tu przed przeczytaniem tego podrozdziału, wpis nabierze sensu później: każda pułapka miała na ciebie zadziałać.

\trapchapter{Rozdział~\ref{ch:models}}{Model to maszyna do przewidywania}
\trapentry{11}{Model to kopia rzeczywistości.}%
  {Model to maszyna do odpowiadania na pytania: liczby na wejściu, liczby na wyjściu, a jedyny uczciwy test to zgodność jego odpowiedzi z tym, co potem zmierzysz.}%
  {\ref{sec:ch01-parts}}
\trapentry{12}{Model, który jest błędny, jest bezużyteczny.}%
  {Każdy model jest błędny, bo coś pomija. Użyteczne pytanie brzmi: czy jest wystarczająco dobry do tego pytania, w tym zakresie.}%
  {\ref{sec:ch01-wrong}}
